% Generated draft: requires mathematical and source review.
\documentclass[11pt]{article}
\usepackage[margin=25mm]{geometry}
\usepackage{fontspec}
\setmainfont{Latin Modern Roman}
\usepackage{amsmath,amssymb,mathtools,unicode-math}
\setmathfont{Latin Modern Math}
\usepackage{enumitem,needspace,xurl,hyperref}
\setlength{\parindent}{0pt}
\setlength{\parskip}{5pt}
\setlength{\emergencystretch}{3em}
\newcommand{\N}{\mathbb{N}}
\newcommand{\Z}{\mathbb{Z}}
\newcommand{\Q}{\mathbb{Q}}
\newcommand{\R}{\mathbb{R}}
\newcommand{\C}{\mathbb{C}}
\newcommand{\F}{\mathbb{F}}
\newcommand{\A}{\mathbb{A}}
\newcommand{\PP}{\mathbb{P}}
\newcommand{\E}{\mathbb{E}}
\newcommand{\eps}{\varepsilon}
\newcommand{\dd}{\,\mathrm{d}}
\begin{document}
\section*{UnsolvedMath ID: 30004368}
\textbf{Source problem code:} catalogue entry
\subsection{Definitions and mathematical statement}
The supplied question is: ``Is every presilting complex with \(n\) indecomposable summands already silting?'' The excerpt does not specify the ambient category, the meaning of \(n\), or how summands are counted. Consequently, the following is an explicitly conditional interpretation requiring confirmation against the source report.

Fix a field \(k\) and a finite-dimensional \(k\)-algebra \(A\). Let \(\mathcal T=K^b(\operatorname{proj} A)\) be the homotopy category of bounded complexes of finitely generated projective right \(A\)-modules, and let \(n\) be the number of isomorphism classes of simple right \(A\)-modules. For \(P\in\mathcal T\), write \(P[i]\) for its shift by \(i\). The complex \(P\) is presilting if
\[
\operatorname{Hom}_{\mathcal T}(P,P[i])=0
\qquad\text{for every integer }i>0.
\]
Let \(\operatorname{thick}_{\mathcal T}(P)\) denote the smallest full triangulated subcategory of \(\mathcal T\) containing \(P\) and closed under direct summands. A presilting complex \(P\) is silting if \(\operatorname{thick}_{\mathcal T}(P)=\mathcal T\). A complex is basic if its indecomposable direct summands are pairwise nonisomorphic.

Under this conditional interpretation, the question is whether every basic presilting complex
\[
P\cong\bigoplus_{j=1}^{n}P_j,
\qquad P_j\text{ indecomposable},\qquad
P_j\not\cong P_\ell\ \text{for }j\ne\ell,
\]
satisfies
\[
\operatorname{thick}_{\mathcal T}(P)=\mathcal T.
\]
The finite-dimensional algebra setting, interpretation of \(n\), basicness convention, and scope of the algebra and field parameters are not established by the supplied excerpt. No bound on the degrees occupied by the complex is stated. These qualifications must be resolved before this conditional formulation can be accepted as the precise source statement.
\subsection{Short English statement}
The question asks whether a presilting complex having the full intended number of indecomposable summands must generate its ambient category and hence be silting. The supplied excerpt leaves the category, the number \(n\), and the summand-counting convention unspecified; the displayed algebraic formulation is conditional and requires human review.
\subsection{Sources}
\begin{itemize}
\item UnsolvedMath local registry, permanent ID \(30004368\), title ``Completion of Presilting Complexes,'' source problem code \texttt{catalogue entry}, in the Oberwolfach Reports Open Problems collection. The record lists proposed year 2020 and supplies no individual proposer.
\item Supplied report citation: \emph{Representation Theory of Quivers and Finite Dimensional Algebras} (2021), Oberwolfach Reports, pp.~143--230, DOI 10.4171/owr/2020/3. \url{https://doi.org/10.4171/owr/2020/3}. The supplied background describes these reports as documenting mathematical workshops and their open-problem discussions.
\item Supplied literature reference: Liu--Zhou, \emph{A negative answer to the complement question}. \url{https://arxiv.org/abs/2302.12502}.
\item Supplied literature reference: Liu--Zhou, \emph{Journal of Algebra}. \url{https://doi.org/10.1016/j.jalgebra.2024.12.020}.
\item Literature assessment checked at \texttt{2026-08-22}, according to the supplied dated curation, not independently verified here: the record labels the question open and states that a 2023--2025 counterexample disproves the stronger complement property for presilting complexes but does not settle the full-number-of-summands question. The supplied background explicitly calls for independent verification of the time-sensitive status and cited claims.
\item The registry labels the statement \texttt{corrected\_verified} and reports a manual semantic, typographic, and LaTeX review claiming self-containedness and preservation of the target. These are supplied review assertions; they do not resolve the contextual omissions identified in this draft.
\end{itemize}
\end{document}
